% !TeX root = ../complete.tex
\section{حل تمرین‌های \lr{Power}}
\label{sec:power-exercises}

چهار تمرین نخست همان تمرین‌های تعیین‌شده در اسلاید ۲۴ هستند.
تمرین تکمیلیِ انتهای بخش، روش تخمین فعالیت و جمع انرژی گره‌ها را
برای مقایسه دو پیاده‌سازی به کار می‌گیرد
\cite[اسلاید ۲۴]{powerslides}\cite[صفحات ۲۰۹ و ۲۱۰]{weste}.

\subsection{تمرین \secnum{۵.۱}، از مساحت به توان تراشه}
\label{sec:power-exercise-chip}

\textbf{خواسته.}
برای تراشه‌ای با منطق تصادفی و \lr{Activity Factor} متوسط \(0.1\)،
توان دینامیکی را تخمین بزنید. ظرفیت سوئیچینگ در فرایند \lr{Standard Cell}
به‌ازای هر میلی‌متر مربع، \qty{450}{\pico\farad} است. مساحت تراشه
\qty{70}{\square\milli\metre}، فرکانس ساعت \qty{450}{\mega\hertz}
و ولتاژ تغذیه \qty{0.9}{\volt} است
\cite[تمرین \secnum{۵.۱}، صفحه ۲۰۹]{weste}.

\textbf{حل.}
ابتدا می‌خواهیم همه ظرفیت‌هایی را که در این مدل می‌توانند شارژ
شوند، معادل یک ظرفیت کنیم. عدد \(c_A\) ظرفیت به‌ازای یک واحد
مساحت است و \(A\) تعداد آن واحدهای مساحت را می‌دهد. بنابراین
در حاصل‌ضرب \(C=c_AA\)، واحد میلی‌متر مربع در صورت و مخرج ساده
می‌شود. برای تبدیل پیکوفاراد به نانوفاراد نیز از این استفاده
می‌کنیم که هر نانوفاراد هزار پیکوفاراد است.
\[
  C=(\qty[per-mode=symbol]{450}{\pico\farad\per\square\milli\metre})
       (\qty{70}{\square\milli\metre})
     =\qty{31500}{\pico\farad}=\qty{31.5}{\nano\farad}.
\]
در رابطه \eqref{eq:power-node-sum}، هر گره با ظرفیت \(C_i\)
و فعالیت \(\alpha_i\) به‌اندازه \(\alpha_i C_i\) به ظرفیت مؤثر
اضافه می‌کند. اگر \(C=\sum_i C_i\) ظرفیت کل باشد، می‌خواهیم
همان ظرفیت مؤثر را با یک فعالیت نماینده \(\bar\alpha_C\) و
حاصل‌ضرب \(\bar\alpha_C C\) بنویسیم. با برابر گرفتن دو بیان
و تقسیم بر \(C>0\)، میانگین فعالیت وزن‌دار با ظرفیت به دست می‌آید.
\begin{equation}
  \bar\alpha_C
   =\frac{\sum_i\alpha_i C_i}{\sum_i C_i}
   =\sum_i\frac{C_i}{C}\,\alpha_i.
  \label{eq:power-capacitance-weighted-activity}
\end{equation}
وزن \(C_i/C\) سهم هر گره از ظرفیت کل است. برای نمونه، فرض کنید
گره اول ظرفیت \qty{1}{\femto\farad} و فعالیت \(1/4\) دارد
و گره دوم با ظرفیت \qty{3}{\femto\farad} ثابت است. میانگین
ساده فعالیت \((1/4+0)/2=1/8\) است. میانگین وزن‌دار از
\((1/4)(1/4)+(3/4)(0)=1/16\) به دست می‌آید. در این مثال
به‌کارگیری میانگین ساده با همان ظرفیت کل و در ولتاژ و فرکانس
یکسان، توان سوئیچینگ را دو برابر برآورد می‌کند.

در برآورد تراشه، مقدار داده‌شده \(0.1\) را نماینده
\(\bar\alpha_C\) می‌گیریم. پس ظرفیت مؤثر
\(\bar\alpha_C C=\qty{3.15}{\nano\farad}\) است. با صرف‌نظر از توان
جریان مستقیمِ گذرا، از رابطه \eqref{eq:power-switching} ابتدا انرژی
متوسطی را که در هر دوره صرف شارژ می‌شود به دست می‌آوریم.
\[
  E_{\mathrm{cycle}}=\bar\alpha_C C V_{DD}^{2}
    =(\qty{3.15}{\nano\farad})(\qty{0.9}{\volt})^2
    =\qty{2.5515}{\nano\joule}.
\]
در این گام، حاصل برحسب ژول به‌ازای هر دوره است. فرکانس تعداد دوره‌ها
در ثانیه است و ضرب آن در انرژی به‌ازای هر دوره، توان متوسط را می‌دهد.
حاصل با گردکردن تا دو رقم اعشار چنین است.
\[
  P_{\mathrm{dyn}}\approx E_{\mathrm{cycle}} f
     =(\qty{2.5515e-9}{\joule})(\qty{450e6}{\hertz})
     \approx\qty{1.15}{\watt}.
\]
ارقام میانی برای نمایش محاسبه حفظ شده‌اند و نتیجه یک برآورد در
مرتبه‌ی یک وات است. با نصف شدن فعالیت و ثابت ماندن ظرفیت، ولتاژ
و فرکانس، مدل نصف شدن این مؤلفه توان را پیش‌بینی می‌کند.

\subsection{تمرین \secnum{۵.۴}، خواندن فعالیت از شکل موج}
\label{sec:power-exercise-waveform}

\textbf{خواسته.}
برای شکل موج تمرین \secnum{۵.۴} کتاب با فرکانس ساعت
\qty{1}{\giga\hertz}، \lr{Activity Factor} را پیدا کنید. در شکل
\ref{fig:power-exercise-waveform}، محور افقی زمان و محور عمودی
نسبت ولتاژ سیگنال \(v\) به ولتاژ تغذیه \(V_{DD}\) است. بنابراین
مقدار یک روی محور عمودی به معنی سطح منطقی بالاست
\cite[تمرین \secnum{۵.۴} و شکل \secnum{۵.۳۴}، صفحه ۲۰۹]{weste}.

\begin{figure}[htbp]
  \centering
  % Shared redraw of the exercise waveform for the notes and submission.
\begin{latin}
\begin{tikzpicture}[x=1.08cm,y=1.30cm]
  \draw[->] (0,-0.30) -- (10.75,-0.30) node[right] {\(t\,(\unit{\nano\second})\)};
  \node[left] at (-0.08,1.24) {\(v/V_{DD}\)};
  \node[left] at (-0.08,0) {\rl{۰}};
  \node[left] at (-0.08,1) {\rl{۱}};
  \foreach \x/\ticklabel in {0/۰,1/۱,2/۲,3/۳,4/۴,5/۵,6/۶,7/۷,8/۸,9/۹,10/۱۰} {
    \draw (\x,-0.25) -- (\x,-0.36) node[below] {\rl{\ticklabel}};
  }
  \draw[gray,densely dashed] (1,-0.24) -- (1,0.50);
  \draw[gray,densely dashed] (4,-0.24) -- (4,0.50);
  \draw[very thick] (0,0) -- (0.75,0) -- (1.25,1) -- (1.75,1)
    -- (2.25,0) -- (3.75,0) -- (4.25,1) -- (7.75,1) -- (8.25,0) -- (10,0);
  \fill (1,0.50) circle (1.5pt) (4,0.50) circle (1.5pt);
\end{tikzpicture}
\end{latin}

  \caption{شکل موج تمرین \secnum{۵.۴} بین صفر تا ده نانوثانیه.
    نقطه‌ها محل تقریبی دو گذار صعودی را نشان می‌دهند. بازترسیم از
    شکل \secnum{۵.۳۴} کتاب \cite{weste}.}
  \label{fig:power-exercise-waveform}
\end{figure}

\textbf{حل.}
برای خواندن نمودار، از چپ به راست در امتداد محور زمان حرکت کنید.
ابتدا سیگنال در سطح صفر است. یک بار نزدیک زمان \qty{1}{\nano\second}
و بار دوم نزدیک زمان \qty{4}{\nano\second} به سطح یک می‌رود. هر بالا رفتن
متناظر با ورود بار به خازن است. گذارهای نزولی در این شمارش سهمی
ندارند، چون انرژیِ ذخیره‌شده را تخلیه می‌کنند. پس صورت کسرِ رابطه
\eqref{eq:power-activity-count} برابر دو است.

برای یافتن مخرج، \(f=\qty{1e9}{\hertz}\) یعنی \(10^9\) دوره در
هر ثانیه. معکوس آن، زمان یک دوره است. از تقسیم طول بازه بر این
دوره، تعداد دوره‌های کامل در بازه به دست می‌آید.
\[
  T_{\mathrm{clk}}=\frac{1}{f}=\qty{1e-9}{\second}
     =\qty{1}{\nano\second},\qquad
  N_{\mathrm{clk}}=\frac{\qty{10}{\nano\second}}{\qty{1}{\nano\second}}=10.
\]
بنابراین \(\alpha=2/10=0.2\) و بی‌بعد است. نرخ شارژ نیز از تقسیم
دو رویداد بر ده نانوثانیه \qty{200}{\mega\hertz} می‌شود. تقسیم
همین نرخ بر فرکانس ساعت دوباره \(0.2\) می‌دهد.

\subsection{تمرین \secnum{۵.۷}، احتمال سطح یک و گذار خروجی}
\label{sec:power-exercise-gates}

\textbf{خواسته.}
احتمال‌های جدول \secnum{۵.۱} کتاب را برای شش گیت آن استخراج
کنید. جدول کتاب احتمال یک‌بودن خروجی را می‌دهد. از آن احتمال،
با مشخص کردن فرض زمانی، فعالیت خروجی را نیز محاسبه می‌کنیم
\cite[جدول \secnum{۵.۱} و تمرین \secnum{۵.۷}، صفحات ۱۸۷ و ۲۰۹]{weste}.

\textbf{حل.}
در گام اول فقط یک دوره را بررسی می‌کنیم. برای ورودی \(X\)،
\(p_X=\Pr(X=1)\) است. چون \(X\) تنها صفر یا یک است، احتمال
صفر بودن آن \(\bar p_X=1-p_X\) می‌شود. استقلال ورودی‌ها در یک
دوره اجازه می‌دهد برای یک ترکیب مشخص، احتمال مقدارِ هر ورودی را
در احتمال مقدارِ ورودی دیگر ضرب کنیم. برای پیدا کردن ترکیب‌های
مطلوب، از \NotesTopic{tab:foundation-gates}{جدول صحت گیت‌ها}
شروع کنید.

\begin{itemize}
  \item برای \lr{AND2}، تنها زوج دارای خروجی یک \((1,1)\) است.
    پس احتمال آن \(p_Ap_B\) می‌شود. \lr{AND3} سه ورودی دارد
    و تنها ترکیبِ موفق \((1,1,1)\) است. با استقلال مشترکِ هر سه
    ورودی، احتمال آن \(p_Ap_Bp_C\) می‌شود.
  \item در \lr{OR2}، یافتن تنها حالتِ دارای خروجی صفر ساده‌تر
    است. این حالت \((0,0)\) است و احتمال \(\bar p_A\bar p_B\)
    دارد. کم کردن آن از یک احتمال خروجی یک را می‌دهد. در
    \lr{NOR2} همین حالت به‌تنهایی خروجی یک دارد.
  \item \lr{NAND2} نتیجه \lr{AND2} را معکوس می‌کند. هر حالتی
    به‌جز \((1,1)\) خروجی یک دارد و احتمال آن \(1-p_Ap_B\) است.
  \item در \lr{XOR2}، خروجی یک است اگر دقیقاً یکی از ورودی‌ها
    یک باشد. حالت \((1,0)\) احتمال \(p_A\bar p_B\) و حالت
    \((0,1)\) احتمال \(\bar p_Ap_B\) دارد. چون هیچ زوجی هر دو
    شرط را هم‌زمان ندارد، حاصل جمع دو احتمال جواب است.
\end{itemize}

در گام دوم دو دوره متوالی را در نظر می‌گیریم. برای یک گذار صعودی،
خروجی اول باید صفر و خروجی بعدی یک باشد. اگر بردار ورودی در دو
دوره مستقل و هم‌توزیع باشد و گذار اضافی میان نمونه‌ها نداشته باشیم،
مقدار خروجی در دو دوره نیز مستقل است. چون خروجی هر دوره فقط تابع
ورودی‌های همان دوره است، می‌توان از رابطه
\eqref{eq:power-activity-independent} استفاده کرد.
برای نمونه، با احتمال \(1/2\) برای هر ورودی، در \lr{OR2}
احتمال صفر بودن خروجی \((1/2)(1/2)=1/4\) و احتمال یک‌بودن آن
\(3/4\) است. بنابراین فعالیت خروجی \((1/4)(3/4)=3/16\) است.
جدول \ref{tab:power-exercise-gates} حاصل همین دو گام را برای هر
شش گیت نشان می‌دهد.

\begin{table}[htbp]
  \centering
  \caption{رابطه احتمال یک‌بودن خروجی و نمونه محاسبه فعالیت.
    دو ستون نمونه از ورودی‌های مستقل با احتمال \(1/2\) و
    استقلال بین دوره‌ها استفاده می‌کنند. روابط احتمال از جدول
    \secnum{۵.۱} کتاب استخراج شده‌اند \cite{weste}.}
  \label{tab:power-exercise-gates}
  \begin{tabular}{rccc}
    \toprule
    گیت & احتمال \(p_Y\) & \(p_Y\) در نمونه & \(\alpha_Y\) در نمونه \\
    \midrule
    \lr{AND2}  & \(p_Ap_B\)                 & \(1/4\) & \(3/16\) \\
    \lr{AND3}  & \(p_Ap_Bp_C\)              & \(1/8\) & \(7/64\) \\
    \lr{OR2}   & \(1-\bar p_A\bar p_B\)    & \(3/4\) & \(3/16\) \\
    \lr{NAND2} & \(1-p_Ap_B\)               & \(3/4\) & \(3/16\) \\
    \lr{NOR2}  & \(\bar p_A\bar p_B\)      & \(1/4\) & \(3/16\) \\
    \lr{XOR2}  & \(p_A\bar p_B+\bar p_Ap_B\) & \(1/2\) & \(1/4\) \\
    \bottomrule
  \end{tabular}
\end{table}

یک راه کنترل این است که خروجی مکمل را بگیریم. با عوض کردن
\(p_Y\) و \(1-p_Y\)، دو عامل در ضرب فعالیت فقط جابه‌جا می‌شوند.
پس با فرض‌های همین مسئله، فعالیت \lr{AND2} و \lr{NAND2} برابر
است. هر دو \lr{OR2} و \lr{NOR2} نیز همین خاصیت را دارند.

\subsection{تمرین \secnum{۵.۱۰}، از قید تأخیر به عرض کلید}
\label{sec:power-exercise-header}

\textbf{خواسته.}
عرض کلید \lr{pMOS Header} در یک مدار \lr{Power Gating} با فرایند
\qty{65}{\nano\metre} را چنان تعیین کنید که افزایش تأخیر کمتر از
۲ درصد بماند. جریان فعال بلوک \qty{100}{\milli\ampere} و
حاصل‌ضرب مقاومت روشن کلید در عرض آن تقریباً
\qty{2.5}{\kilo\ohm\micro\metre} است
\cite[تمرین \secnum{۵.۱۰}، صفحه ۲۱۰]{weste}.

\textbf{حل.}
شکل \ref{fig:power-gating-header} نشان می‌دهد که جریان تغذیه‌ی بلوک
در حالت فعال از کلید می‌گذرد. مقاومت کلید افت ولتاژ ایجاد می‌کند
و بلوک عملاً با \(V_{DDV}=V_{DD}-\Delta V\) کار می‌کند. هرچه
\(\Delta V\) بیشتر باشد، طبق مدل بخش \ref{sec:power-header-delay}
بلوک کندتر می‌شود. ابتدا از محدودیت تأخیر به محدودیت ولتاژ می‌رسیم.

برای عددگذاری، \(V_{DD}=\qty{1}{\volt}\) را از مثال ۵.۵ کتاب
می‌گیریم. آن مثال برای افت کوچک، درصد افزایش تأخیر را تقریباً
برابر درصد افت تغذیه می‌گیرد. رابطه
\eqref{eq:power-droop-delay-sensitivity} نشان می‌دهد که این ضریب
به آستانه و مدل جریان وابسته است. با همان تقریبِ مثال کتاب، برای
کمتر از ۲ درصد افزایش تأخیر لازم است
\(\Delta V/V_{DD}<2/100\) باشد. پس حد افت از ضرب \(0.02\) در
یک ولت به \qty{20}{\milli\volt} می‌رسد
\cite[مثال ۵.۵، صفحات ۱۹۸ و ۱۹۹]{weste}.

اکنون قانون اهم را روی کلید اعمال می‌کنیم. از رابطه
\eqref{eq:power-header-droop}، تقسیم حد افت بر جریان، مرز مقاومت
را می‌دهد. نسبت میلی‌ولت به میلی‌آمپر برابر ولت بر آمپر است و به
اهم تبدیل می‌شود.
\[
  R_h < \frac{\qty{20}{\milli\volt}}{\qty{100}{\milli\ampere}}
       =\qty{0.20}{\ohm}.
\]
اگر \(r_w=\qty{2.5e3}{\ohm\micro\metre}\) و \(W\) عرض مؤثر
کلید باشد، \(R_h\approx r_w/W\) است. همه کمیت‌ها مثبت‌اند. برای
کوچک‌تر شدن مقاومت، مخرج باید بزرگ‌تر شود. با ضرب نامساوی در
\(W\) و تقسیم بر حد مقاومت، داریم
\[
  W>\frac{\qty{2.5e3}{\ohm\micro\metre}}{\qty{0.20}{\ohm}}
    =\qty{12.5e3}{\micro\metre}=\qty{12.5}{\milli\metre}.
\]
در مرز \qty{12.5}{\milli\metre}، افت مدل دقیقاً به
\qty{20}{\milli\volt} می‌رسد. با شرط «کمتر از» در صورت مسئله،
عرض مؤثر باید از این مرز بیشتر باشد. اگر کلید با
ترانزیستورهای موازی ساخته شود، منظور مجموع عرض‌هاست. مقدار مرزی
یک تخمین مدل است و برای انتخاب عرض در یک مدار مشخص، بایاس کلید و
وابستگی تأخیر به ولتاژ آن بلوک لازم‌اند.

\subsection{تمرین تکمیلی، \lr{AND} چهارورودی با ساختار درختی و زنجیره‌ای}
\label{sec:power-exercise-and-topology}

\textbf{خواسته.}
با سه گیت \lr{AND2}، تابعی بسازید که فقط با یک‌بودن ورودی‌های
\(A,B,C,D\) یک شود. در طرح درختی، ابتدا \(M=A\land B\) و
\(N=C\land D\) و سپس \(Y=M\land N\) محاسبه می‌شوند. در طرح
زنجیره‌ای، پس از \(M=A\land B\)، گره \(Q=M\land C\) و سپس
\(Y=Q\land D\) به دست می‌آید. اگر همه‌ی ورودی‌ها درون هر دوره و
بین دوره‌ها مستقل و هم‌احتمال باشند، کدام طرح انرژی سوئیچینگ
کمتری دارد؟ ظرفیت هر گره خروجی گیت را \(C_L\)، تغذیه را در هر
دو طرح یکسان و \lr{Glitch} را ناچیز بگیرید. ایده مقایسه از مثال
۵.۲ کتاب گرفته شده و پیاده‌سازی در اینجا با گیت‌های \lr{AND2}
در سطح منطقی ساده شده است. در مقایسه با
\NotesTopic{sec:power-slide-and-activity}{مدار اسلاید ۱۶}، به‌کارگیری
\lr{AND2} برای دو گره میانی درخت، احتمال یک بودن هر گره را
از \(3/4\) به \(1/4\) تغییر می‌دهد و فعالیت هر گره همچنان
\(3/16\) است \cite[صفحات ۱۸۷ و ۱۸۸]{weste}.

\textbf{حل.}
در هر دو طرح، \(M\) تنها زمانی یک می‌شود که دو ورودی نخست یک
باشند. بنابراین احتمال آن \(1/4\) و فعالیت آن \(3/16\) است.
در طرح درختی، دو ورودیِ \(N\) نیز احتمال‌های یکسانی دارند، پس
فعالیت \(N\) برابر \(3/16\) است. در طرح زنجیره‌ای، \(Q\) با
یک‌بودن سه ورودی مستقل به یک می‌رسد. در نتیجه، احتمال آن
\((1/2)^3=1/8\) و فعالیت آن \((1/8)(7/8)=7/64\) می‌شود.
مقدار \(Y\) در هر دو پیاده‌سازی تنها در یک حالت از شانزده حالت
هم‌احتمال یک است. پس \(p_Y=1/16\) و \(\alpha_Y=15/256\) است.

برای به دست آوردن انرژی سوئیچینگ متوسط در هر دوره، از مجموع ظرفیت
وزن‌دار در رابطه \eqref{eq:power-node-sum} استفاده می‌کنیم. با
ظرفیت‌های برابر، برای مقایسه کافی است مجموع فعالیت خروجی گیت‌ها
را در هر طرح پیدا کنیم. با تبدیل کسرها به مخرج مشترک داریم
\[
  \begin{aligned}
    \alpha_M+\alpha_N+\alpha_Y
      &=\frac{3}{16}+\frac{3}{16}+\frac{15}{256}=\frac{111}{256},\\
    \alpha_M+\alpha_Q+\alpha_Y
      &=\frac{3}{16}+\frac{7}{64}+\frac{15}{256}=\frac{91}{256}.
  \end{aligned}
\]
پس انرژی طرح زنجیره‌ای در این مدل \(91/111\) انرژی طرح درختی است.
صرفه‌جویی نسبت به طرح درختی \((111-91)/111\)، یعنی حدود ۱۸ درصد
است. طرح زنجیره‌ای در مسیر برخی ورودی‌ها سه مرحله دارد و مسیرهای
طرح درختی دو مرحله دارند. اگر مهلت زمانی، ظرفیت‌های نابرابر یا
\lr{Glitch} در نتیجه اثرگذار باشند، باید انرژی و تأخیر را برای
بار و شکل موج واقعی هر طرح دوباره مقایسه کرد.
